\chapter{Marco Metodológico}
% BORRADOR para editar a tu voz (estilo impersonal UAX). Recomendado: 5-15 páginas.

En este capítulo se describe el método seguido para desarrollar y evaluar el trabajo. Se
detallan la población y la muestra empleadas, los objetivos de investigación y sus hipótesis,
las variables medidas y los instrumentos de evaluación, el diseño experimental y, finalmente,
el modelo de análisis de datos. El enfoque combina el desarrollo de un sistema de generación de
informes con un marco de evaluación de la explicabilidad, y su validación tanto sobre un modelo
propio como sobre un modelo multimodal del estado del arte.

\section{Población y muestra}
La población objeto de estudio son los estudios de radiografía de tórax (CXR) acompañados de su
informe clínico. Se han utilizado dos conjuntos de datos de acceso libre, seleccionados por su
disponibilidad sin restricciones frente a alternativas como MIMIC-CXR, cuyo acceso exige
credenciales de PhysioNet no disponibles para este trabajo.

El conjunto principal es \textbf{Open-I / Indiana University CXR} \parencite{demnerfushman2016},
que reúne 3.955 estudios y 7.470 imágenes (vistas frontal y lateral) con informes narrativos
estructurados en secciones \textit{FINDINGS} e \textit{IMPRESSION}. La partición en
entrenamiento, validación y prueba (70/15/15) se realiza \textbf{por identificador de estudio},
de modo que ambas vistas de un mismo paciente permanecen en la misma partición y se evita la
fuga de información entre conjuntos.

El segundo conjunto es \textbf{NIH ChestX-ray8} \parencite{wang2017}, del que se emplean las 984
anotaciones con \textbf{caja delimitadora (bounding box)} correspondientes a ocho patologías
(atelectasia, cardiomegalia, derrame, infiltrado, masa, nódulo, neumonía y neumotórax). Estas
cajas constituyen la \textbf{referencia objetiva de localización} que permite evaluar el anclaje
clínico de la atención sin depender de un juicio subjetivo. La restricción de datos narrativos
—Open-I es prácticamente el único conjunto libre con informes completos— motiva que la mejora
del modelo propio se aborde mediante \textbf{destilación de conocimiento} desde modelos
preentrenados libres, convirtiendo la limitación en una decisión metodológica.

El tercer conjunto es \textbf{PadChest-GR} \parencite{padchestgr2024}, subconjunto anclado del corpus
español PadChest, del que se emplea una muestra de estudios con \textbf{hallazgos presentes y ausentes
anotados por radiólogo} y sus cajas. PadChest-GR cumple una doble función: proporciona una
\textbf{segunda fuente de verdad objetiva} —independiente del juez basado en modelo de lenguaje— y
constituye un \textbf{dominio distinto} (población española) respecto a Open-I (Estados Unidos), lo
que permite evaluar la generalización \emph{fuera de distribución}. El cruce de ambos dominios con
dos generadores da lugar al diseño \(2\times2\) que vertebra la validación.

En cuanto al \textbf{detalle y el preprocesado}, Open-I reúne 3.955 estudios y 7.470 imágenes
(proyecciones frontal y lateral); su partición 70/15/15 se realiza por identificador de estudio para
que ambas vistas de un paciente permanezcan juntas y no se produzca fuga de información. De NIH se
emplean las 984 anotaciones con caja de ocho patologías. Del subconjunto de PadChest-GR analizado
—382 estudios, de los cuales 120 son normales y 262 presentan algún hallazgo— se extraen los hallazgos
presentes y ausentes anotados por radiólogo. En todos los casos, las imágenes se convierten a escala de
grises de tres canales, se redimensionan a la resolución de entrada del codificador correspondiente
(224 píxeles para el \textit{Vision Transformer}, 518 para RAD-DINO) y se normalizan con las
estadísticas del preentrenamiento. El texto de referencia se limpia de secciones administrativas,
conservando \textit{Findings} e \textit{Impression}.

\section{Objetivos de investigación}
La investigación empírica se organiza en torno a tres hipótesis:
\begin{itemize}
  \item \textbf{H1.} La atención de un modelo multimodal de CXR está \textit{anclada
  clínicamente}, es decir, se concentra sobre la región anatómica de la lesión.
  \item \textbf{H2.} La \textit{forma} del mapa de calor (su morfología: concentración,
  dispersión, etc.) permite detectar cuándo el modelo se equivoca (alucinación o error).
  \item \textbf{H3.} La \textit{incertidumbre generativa} del modelo (las probabilidades de su
  propia salida) permite detectar dichos errores.
  \item \textbf{H4.} Dicha señal de incertidumbre es \textit{agnóstica al modelo y robusta al
  dominio}: se transfiere a otro generador y a otra población distinta de la de entrenamiento.
\end{itemize}
Como se mostrará, H1 se confirma y H3 se sostiene, mientras que H2 se refuta de forma robusta: la
fiabilidad no reside en la morfología de la explicación sino en la incertidumbre de la decisión del
modelo. H4 se confirma con un matiz relevante: la señal se transfiere entre modelos y se mantiene
fuera de distribución para un generador de calidad suficiente, mientras que un generador débil la
pierde fuera de distribución al recaer en un informe-plantilla emitido con confianza.

\section{Variables intervinientes e instrumentos de evaluación}
Las \textbf{variables} evaluadas se agrupan en tres familias. (i) \textit{Morfología del mapa de
calor}: concentración, dispersión (entropía, índice de Gini), agrupamiento y autocorrelación
espacial. (ii) \textit{Incertidumbre generativa}: probabilidad del token menos seguro
(\texttt{min\_prob}) y entropía predictiva por token. Una decisión de diseño relevante es emplear la
\textbf{probabilidad mínima de \emph{contenido}} —excluyendo el token de apertura del informe, cuya
incertidumbre es de estilo y no de contenido—, que mejora sustancialmente la señal (véase
Cap.~5). (iii) \textit{Dependencia de la imagen y consistencia}: fracción de atención dirigida a los
tokens visuales frente al texto, persistencia del hallazgo ante una imagen en blanco
(\emph{blank\_sim}) y variabilidad entre muestras (auto-consistencia tipo SelfCheckGPT). La
\emph{blank\_sim} resulta además un indicador de la calidad del generador: cuanto más depende de la
imagen (menor similitud con la salida en blanco), mejor anclado está el modelo.

Los \textbf{instrumentos} de evaluación son: las métricas de anclaje \textit{pointing-game},
\textit{energy-in-bbox} y \textit{lift} (energía dentro de la caja dividida por su área, es decir,
la mejora sobre el azar); las métricas de discriminación (área bajo la curva ROC, AUC) con
\textit{test de permutación} e \textbf{intervalos de confianza por \emph{bootstrap}}; y el
etiquetado del error, para el que se emplean \textbf{tres fuentes de verdad independientes} a fin de
evitar la circularidad de un único juez: (a) el \textbf{diagnóstico estructurado oficial} (términos
MeSH de Open-I; hallazgos presentes/ausentes anotados por radiólogo en PadChest-GR), objetivo y
reproducible; y (b) el \textbf{juez basado en un modelo de lenguaje} (metodología tipo G-Eval), que
compara la \textbf{narrativa completa} generada con la referencia y la clasifica como correcta,
alucinada u omitida según una rúbrica predefinida. La concordancia entre las tres fuentes (por
ejemplo, AUC 0,87 con juez, 0,82 con MeSH oficial sobre Open-I) sustenta la fiabilidad del
etiquetado. El juez se ejecuta con escritura incremental y de forma reproducible sobre cada informe.

\begin{table}[t]\centering
\caption{Conjuntos de datos empleados y su función en el trabajo.}
\label{tab:datos}
\begin{tabular}{llll}
\toprule
Conjunto & Contenido & Función & Verdad \\ \midrule
Open-I & Informes + MeSH & Estudio principal (in-dist.) & MeSH + juez \\
NIH ChestX-ray8 & Cajas (8 patologías) & Anclaje objetivo & Cajas \\
PadChest-GR & Presentes/ausentes + cajas & Réplica (dominio, OOD) & Radiólogo + juez \\ \bottomrule
\end{tabular}
\end{table}

\section{El marco de explicabilidad}
El marco de explicabilidad implementado evalúa dos ejes complementarios: \emph{dónde} mira el modelo y
si \emph{de verdad} depende de ello. El \textbf{mapa de calor} se obtiene agregando la atención que
los tokens del informe dirigen a los tokens visuales del codificador. Para los modelos de tipo LLaVa,
los tokens visuales forman una rejilla de $37\times37$ (1369 parches); se promedia la atención sobre
capas y cabezas, se restringe a los tokens visuales y se reordena en la rejilla, obteniendo un mapa
espacial que se superpone a la imagen. Se descarta la contribución de los tokens ``sumidero''
(posiciones de atención residual sin significado espacial). Sobre este mapa se calculan las
\textbf{métricas de anclaje} frente a las cajas reales de NIH: el \textit{energy-in-bbox} (fracción de
energía de atención dentro de la caja), el \textit{lift} (dicha fracción normalizada por el área,
es decir, la mejora sobre el azar) y el \textit{pointing-game} (si el máximo del mapa cae dentro de la
caja).

La \textbf{fidelidad} de la explicación se evalúa por \textbf{perturbación}: se ocluye la región de
máxima atención y se mide la caída de la log-probabilidad del hallazgo (\emph{occ\_drop}). Una
explicación fiel debería producir una caída apreciable —si la región atendida realmente sustenta la
afirmación—; su ausencia indicaría que la atención no es causalmente responsable de la salida. Como se
detalla en el capítulo de resultados, esta y otras métricas de forma \emph{no} logran predecir la
fiabilidad, lo que motiva el giro hacia la incertidumbre de la decisión.

\section{Arquitectura del generador propio}
El generador propio es un modelo \textbf{codificador--atención cruzada--decodificador} para la
generación de informes de tórax. El \textbf{codificador visual} es un \textit{Vision Transformer}
(ViT) preentrenado que produce, para cada imagen, un embedding global (CLS) y una rejilla de
características por parche; en las variantes más avanzadas se sustituye por \textbf{RAD-DINO}
\parencite{perezgarcia2024} congelado (518\,px, rejilla de $37\times37$ parches, 768 dimensiones),
aprovechando un codificador de dominio de mayor calidad. El \textbf{decodificador} es un modelo de
lenguaje GPT-2 al que se \textbf{inyecta atención cruzada} hacia las características visuales en cada
bloque. Un aspecto clave de diseño es que dicha atención cruzada se introduce mediante una
\textbf{compuerta inicializada a cero}: al comienzo del entrenamiento la contribución visual es nula y
el decodificador se comporta exactamente como el GPT-2 preentrenado, lo que \emph{preserva} sus pesos
y evita el colapso de la generación; la compuerta aprende progresivamente cuánta información visual
incorporar.

El entrenamiento se organiza en \textbf{dos etapas}. En la \textbf{Etapa 1 (alineamiento
visual--textual)} se entrena únicamente el codificador y unas cabezas de proyección con una pérdida
contrastiva de tipo InfoNCE (CLIP): los embeddings de imágenes e informes emparejados deben acercarse
frente a los no emparejados. En la \textbf{Etapa 2 (ajuste \mbox{generativo})} se entrena el modelo completo
con una pérdida de modelado de lenguaje (entropía cruzada sobre el siguiente token). Opcionalmente,
mediante ajuste eficiente en parámetros (PEFT), se congela el GPT-2 base y se entrenan solo las capas
de atención cruzada, reduciendo los parámetros entrenables de unos 200 a unos 5 millones y mitigando
la ceguera a la imagen.

Ante la restricción de datos narrativos, el modelo incorpora \textbf{destilación de conocimiento}, que
se estudia mediante una ablación progresiva (M0--M3): \textbf{M0}, línea base sin destilación; en el
eje textual, se destila desde \textbf{CXR-BERT} \parencite{boecking2022} —bien como torre de texto
congelada que aporta un embedding clínico, bien como objetivo de destilación que el modelo
\emph{student} aprende a predecir—; y en el eje visual (\textbf{M3}), se destila el embedding CLS del
ViT hacia el de RAD-DINO. Esta ablación permite aislar la contribución de cada fuente de conocimiento
externo a la calidad y al anclaje del modelo resultante. La Figura~\ref{fig:arch} esquematiza la
arquitectura y sus dos etapas de entrenamiento.

\begin{figure}[t]\centering
\begin{tikzpicture}[
  node distance=7mm and 20mm,
  enc/.style={draw, rounded corners, align=center, minimum height=16mm, minimum width=22mm, font=\footnotesize, fill=blue!6},
  dec/.style={draw, rounded corners, align=center, minimum height=16mm, minimum width=26mm, font=\footnotesize, fill=orange!8},
  io/.style={draw, align=center, minimum height=9mm, font=\footnotesize, fill=gray!8},
  flow/.style={-{Latex[length=2mm]}, thick}]
  \node[io] (img) {Radiografía};
  \node[enc, right=of img] (enc) {Codificador visual\\(ViT / RAD-DINO)};
  \node[dec, right=of enc] (dec) {Decodificador GPT-2\\+ atención cruzada\\(compuerta $\to 0$)};
  \node[io, right=of dec] (rep) {Informe};
  \draw[flow] (img) -- (enc);
  \draw[flow] (enc) -- node[above=1.5mm, font=\scriptsize, align=center]{parches\\$37{\times}37$} (dec);
  \draw[flow] (dec) -- (rep);
  \node[font=\scriptsize, below=3mm of enc, text width=30mm, align=center] {\textbf{Etapa 1}: alineamiento contrastivo (InfoNCE)};
  \node[font=\scriptsize, below=3mm of dec, text width=34mm, align=center] {\textbf{Etapa 2}: ajuste generativo (modelado de lenguaje)};
\end{tikzpicture}
\caption{Arquitectura del generador propio: codificador visual, decodificador con atención cruzada
mediante compuerta inicializada a cero, y las dos etapas de entrenamiento (alineamiento y generación).}
\label{fig:arch}
\end{figure}

\section{Definición formal de las señales}
Para garantizar la reproducibilidad, se definen formalmente las señales empleadas. Sea un informe
generado la secuencia de tokens $(w_1,\dots,w_n)$ con probabilidades $p_t = P(w_t \mid w_{<t},
\text{imagen})$ asignadas por el modelo en cada paso $t$.

\paragraph{Incertidumbre generativa.} La señal principal es la \textbf{probabilidad mínima de
contenido}, que excluye el token de apertura (cuya incertidumbre es de estilo y no de contenido):
\begin{equation}
u_{\min} \;=\; -\!\!\min_{2 \le t \le n} p_t .
\end{equation}
Se calcula también la \textbf{entropía predictiva} por token, $H_t = -\sum_{v} p_t(v)\log p_t(v)$, de
la que se toma el máximo $H_{\max}=\max_t H_t$. Para eliminar su confusión con la longitud, la
entropía se \textbf{residualiza} respecto al número de tokens $n$ mediante una regresión lineal
$H_{\max} = \alpha + \beta n + \varepsilon$, tomando el residuo $\varepsilon$ como señal limpia.

\paragraph{Auto-consistencia.} Dadas $K$ muestras $\{s_1,\dots,s_K\}$ generadas por muestreo
estocástico, se mide su acuerdo con la generación voraz $g$ mediante la \textbf{similitud de Jaccard
de $n$-gramas de carácter}:
\begin{equation}
\mathrm{sim}(a,b) = \frac{|N(a)\cap N(b)|}{|N(a)\cup N(b)|},\qquad
c = \frac{1}{K}\sum_{k=1}^{K}\mathrm{sim}(g,s_k),
\end{equation}
donde $N(\cdot)$ es el conjunto de $4$-gramas de carácter. Una baja consistencia $c$ indica un
contenido inestable entre muestras, señal candidata de alucinación.

\paragraph{Dependencia de la imagen.} Se genera además un informe con la \textbf{imagen en blanco}
$g_{\varnothing}$ y se calcula $\mathrm{blank\_sim} = \mathrm{sim}(g, g_{\varnothing})$: cuanto menor
es, más depende de la imagen la salida real.

\paragraph{Anclaje objetivo.} Sobre NIH, dado el mapa de calor normalizado $h$ y la caja $B$ de área
relativa $|B|$, el \textbf{lift} mide la mejora sobre el azar de la energía de atención contenida en
la caja:
\begin{equation}
\mathrm{lift} = \frac{\sum_{(i,j)\in B} h_{ij}}{|B|}.
\end{equation}
Un valor $\mathrm{lift}>1$ indica que la atención se concentra en la lesión por encima de lo
esperable por azar.

\section{Protocolo de etiquetado del error}
La variable dependiente —\emph{correcto}, \emph{alucinación} u \emph{omisión}— se determina por dos
mecanismos complementarios. El \textbf{etiquetado objetivo} compara los hallazgos agudos afirmados por
el modelo con la verdad estructurada (MeSH en Open-I; presentes/ausentes en PadChest-GR): se cuenta
alucinación cuando el modelo afirma un hallazgo agudo ausente de la referencia, y omisión cuando
omite uno presente. El \textbf{etiquetado por juez} (G-Eval) somete cada par (informe generado,
referencia) a un modelo de lenguaje que aplica una \textbf{rúbrica fija}: se consideran significativos
únicamente los hallazgos \emph{agudos} (cardiomegalia, derrame, consolidación, edema, neumotórax,
masa, nódulo, neumonía, infiltrado, atelectasia), ignorándose los crónicos, óseos, vasculares,
técnicos o triviales; un informe es \emph{correcto} si coincide con la referencia en los hallazgos
significativos, \emph{alucinación} si afirma uno ausente y \emph{omisión} si pierde uno presente. La
concordancia entre ambos mecanismos valida el etiquetado y mitiga la circularidad de depender de un
único criterio.

\section{Requisitos del sistema}
Aunque el trabajo es fundamentalmente de investigación, el sistema resultante se especifica mediante
requisitos, lo que ordena su diseño y su evaluación. Los \textbf{requisitos funcionales} (RF) son:
\textbf{RF1}, generar el informe de una radiografía de tórax; \textbf{RF2}, extraer en inferencia
las señales de incertidumbre sin generaciones adicionales obligatorias; \textbf{RF3}, calcular el
anclaje de la explicación (mapa de calor o cajas); \textbf{RF4}, estimar y \textbf{calibrar} una
confianza por informe; \textbf{RF5}, presentar informe, anclaje y confianza de forma integrada; y
\textbf{RF6}, priorizar la revisión mediante bandas de riesgo. Los \textbf{requisitos no funcionales}
(RNF) son: \textbf{RNF1}, ejecutarse en una única tarjeta de 12\,GB mediante cuantización;
\textbf{RNF2}, que la señal de confianza sea \textbf{agnóstica al modelo} y no incurra en coste de
inferencia adicional; \textbf{RNF3}, reproducibilidad (extracción resumible, cálculo determinista);
\textbf{RNF4}, interpretabilidad y honestidad de la salida (comunicar las limitaciones); y
\textbf{RNF5}, portabilidad y bajo coste de despliegue.

\section{Arquitectura de la solución}
La solución se organiza en \textbf{cinco capas}. La \textbf{capa de datos} gestiona los conjuntos
(Open-I, NIH, PadChest-GR), su descarga, anonimización y preprocesado. La \textbf{capa de modelo e
inferencia} carga los generadores cuantizados y produce, por imagen, el informe, las probabilidades y
entropías por token, las muestras estocásticas y la atención visual. La \textbf{capa de análisis y
explicabilidad} calcula a partir de lo anterior las señales de incertidumbre, las métricas de anclaje
y las etiquetas de error. La \textbf{capa de calibración} transforma la señal en una probabilidad con
bandas de riesgo (el semáforo). Y la \textbf{capa de presentación} materializa el sistema en la
demostración interactiva. Esta separación por capas favorece la reutilización —el marco de análisis es
\emph{agnóstico al generador}— y facilita sustituir un modelo por otro sin alterar el resto. La
Figura~\ref{fig:pipeline} representa el flujo completo del sistema.

\begin{figure}[t]\centering
\begin{tikzpicture}[
  node distance=6mm and 9mm,
  box/.style={draw, rounded corners, align=center, minimum height=9mm, minimum width=20mm, font=\footnotesize, fill=blue!5},
  data/.style={draw, align=center, minimum height=9mm, font=\footnotesize, fill=gray!8},
  outbox/.style={draw, rounded corners, align=center, minimum height=9mm, font=\footnotesize, fill=green!8},
  arr/.style={-{Latex[length=2mm]}, thick}]
  \node[data] (img) {Radiografía};
  \node[box, right=of img] (gen) {Generador\\(LLaVA-Rad /\\MAIRA-2)};
  \node[data, right=of gen] (sig) {Informe\\Probabilidades\\Atención};
  \node[box, right=of sig] (an) {Análisis:\\señales +\\anclaje};
  \node[box, right=of an] (cal) {Calibración\\(semáforo)};
  \node[outbox, below=8mm of cal] (demo) {Demo: informe\\+ anclaje\\+ confianza};
  \draw[arr] (img) -- (gen);
  \draw[arr] (gen) -- (sig);
  \draw[arr] (sig) -- (an);
  \draw[arr] (an) -- (cal);
  \draw[arr] (cal) -- (demo);
  \node[data, below=8mm of an] (truth) {Verdad:\\MeSH / radiólogo\\/ cajas NIH};
  \draw[arr] (truth) -- (an);
\end{tikzpicture}
\caption{Flujo del sistema: de la radiografía al informe con anclaje y confianza calibrada. El bloque
de análisis se valida contra la verdad objetiva (MeSH, hallazgos de radiólogo, cajas de NIH).}
\label{fig:pipeline}
\end{figure}

\section{Casos de uso}
Se contemplan cuatro casos de uso principales. \textbf{CU-01 (radiólogo): obtener informe y confianza}
—el usuario carga una radiografía y el sistema devuelve el informe, el anclaje y la banda de
confianza—. \textbf{CU-02 (radiólogo): consultar el anclaje} —inspeccionar dónde sustenta el modelo su
afirmación, como apoyo a la lectura—. \textbf{CU-03 (radiólogo): priorizar la revisión} —ordenar la
carga de trabajo atendiendo primero a los informes en zona roja—. Y \textbf{CU-04 (investigador):
auditar y reproducir} —extraer las señales sobre un conjunto y reproducir las métricas de anclaje y
detección—.

\section{Planificación y desarrollo}
El proyecto se ha desarrollado de forma incremental, con control de versiones (Git) y gestión de
tareas por incrementos, organizándose en seis fases que se describen a continuación. Cada fase produce
un artefacto verificable —un componente de software o un resultado experimental— antes de avanzar a la
siguiente, lo que permite congelar resultados y retomar procesos largos entre sesiones. El trabajo de
campo experimental (extracción de señales y etiquetado sobre miles de imágenes) se ejecuta en procesos
reanudables que sobreviven a interrupciones, dado el tiempo de cómputo en hardware de consumo.

\section{Diseño experimental}
El trabajo se estructura en seis fases experimentales.
\begin{enumerate}
  \item \textbf{Modelo propio (POC).} Un codificador visual (Vision Transformer) y un decodificador
  de lenguaje (GPT-2) al que se inyecta atención cruzada en cada bloque mediante una compuerta
  inicializada a cero, que preserva los pesos preentrenados. El entrenamiento sigue dos etapas:
  alineamiento contrastivo (InfoNCE) y ajuste generativo. Se realiza una ablación (M0--M3) que
  introduce progresivamente la destilación de conocimiento textual (CXR-BERT
  \parencite{boecking2022}) y visual (RAD-DINO \parencite{perezgarcia2024}).
  \item \textbf{Elevación a un modelo del estado del arte.} El marco de explicabilidad se aplica a
  LLaVA-Rad, ejecutado en el hardware local mediante cuantización a 8 bits, adaptando la
  extracción de mapas a su arquitectura de tokens visuales antepuestos.
  \item \textbf{Anclaje objetivo.} Validación del anclaje contra las cajas de NIH.
  \item \textbf{Detección de errores.} Estudio multi-señal de la fiabilidad, combinando morfología,
  incertidumbre, dependencia de imagen y consistencia, con etiqueta objetiva (MeSH) y juez.
  \item \textbf{Réplica cruzada (diseño \(2\times2\)).} La señal se valida sobre un \textbf{segundo
  generador del estado del arte, MAIRA-2} \parencite{maira2}, y sobre un \textbf{segundo dominio,
  PadChest-GR}, cruzando ambos factores (dos modelos \(\times\) dos dominios). MAIRA-2 se ejecuta
  también cuantizado a 8 bits; al haber sido entrenado con PadChest, su evaluación
  \emph{in-distribution} se realiza sobre PadChest y la evaluación \emph{fuera de distribución} sobre
  Open-I, ajeno a su entrenamiento. Se aprovecha además su capacidad de \textbf{anclaje explícito
  (cajas)} para contrastar por tercera vía si la localización de la explicación predice la fiabilidad.
  \item \textbf{Calibración y confianza operativa.} La señal final se calibra para expresar la
  fiabilidad como una probabilidad con bandas de riesgo (``semáforo''), y se materializa en una
  \textbf{demostración interactiva} que integra informe, anclaje y confianza.
\end{enumerate}

\section{Métricas de evaluación}
La evaluación cuantitativa se apoya en un conjunto de métricas que conviene definir formalmente.

\paragraph{Discriminación (AUC--ROC).} La capacidad de una señal para separar dos clases se mide con el
\textbf{área bajo la curva ROC}. La curva ROC representa la sensibilidad (tasa de verdaderos positivos)
frente a la tasa de falsos positivos al variar el umbral de decisión. El AUC admite una interpretación
probabilística: es la probabilidad de que una instancia positiva reciba una puntuación mayor que una
negativa elegida al azar, y equivale al estadístico de Mann--Whitney,
\begin{equation}
\mathrm{AUC} = \frac{1}{n_{+}n_{-}}\sum_{i\in\mathcal{P}}\sum_{j\in\mathcal{N}} \mathbf{1}[s_i > s_j],
\end{equation}
donde $\mathcal{P}$ y $\mathcal{N}$ son los conjuntos de positivos y negativos y $s$ las puntuaciones.
Un valor de $0{,}5$ indica ausencia de discriminación y $1{,}0$ discriminación perfecta. Se calcula
mediante rangos para garantizar reproducibilidad.

\paragraph{Significancia (test de permutación).} Para contrastar si un AUC observado supera al azar
—especialmente con clases desbalanceadas— se emplea un \textbf{test de permutación}: se permutan
aleatoriamente las etiquetas un número elevado de veces, se recalcula el AUC en cada permutación y se
obtiene el $p$-valor como la fracción de valores nulos que igualan o superan al observado.

\paragraph{Intervalos de confianza (\textit{bootstrap}).} La incertidumbre de la estimación se cuantifica
por \textbf{\textit{bootstrap} estratificado}: se remuestrean con reemplazo, por separado, los
positivos y los negativos un número elevado de veces (2000 en este trabajo), y el intervalo de
confianza al 95\,\% se toma de los percentiles $2{,}5$ y $97{,}5$ de la distribución de AUC
remuestreados.

\paragraph{Calibración.} Que una señal discrimine no implica que esté \textbf{calibrada}. La
calibración se evalúa con la \textbf{curva de fiabilidad}, que compara la probabilidad predicha con la
tasa de acierto observada por grupos, y con el \textbf{\textit{Brier score}}, el error cuadrático medio
entre la probabilidad predicha $\hat{p}_i$ y el resultado $y_i$:
$\mathrm{Brier}=\frac{1}{n}\sum_i (\hat{p}_i - y_i)^2$. Sobre esta base se definen las bandas de riesgo
del semáforo.

\paragraph{Anclaje y similitud.} El anclaje se mide con el \textit{lift}, el \textit{energy-in-bbox} y
el \textit{pointing-game} ya definidos; la consistencia textual, con la similitud de Jaccard de
$n$-gramas de carácter.

\section{Método científico y control de confusores}
El trabajo sigue un enfoque \textbf{hipotético-deductivo}: se formulan hipótesis falsables (H1--H4,
Sección~4.2) y se diseñan experimentos que puedan refutarlas. Tres principios metodológicos guían la
investigación. Primero, la \textbf{validación contra verdad objetiva}: siempre que es posible, las
señales se contrastan con referencias externas (cajas reales, diagnóstico estructurado) en lugar de con
otras señales del propio modelo, y el etiquetado del error se \textbf{triangula} con tres fuentes
independientes. Segundo, el \textbf{control de confusores}: se identifican y neutralizan las variables
espurias que podrían simular una señal —la longitud del informe (por residualización), la posición del
token (usando la probabilidad de contenido) o el tipo de mapa de calor (mediante comparaciones
emparejadas)—. Tercero, la \textbf{replicación}: los hallazgos se comprueban en un diseño cruzado de dos
modelos por dos dominios antes de darlos por robustos. Este rigor resultó decisivo: varias señales de
apariencia prometedora (el máximo de atención, la entropía cruda, la relación \emph{lift}--acierto) se
revelaron \textbf{artefactos} que solo un control cuidadoso de confusores permitió descartar,
subrayando la importancia del método sobre el resultado aparente.

\section{Modelo de análisis de datos}
El análisis es fundamentalmente cuantitativo. La capacidad discriminativa de cada señal se mide
mediante AUC univariante; las señales se combinan con regresión logística evaluada por
\textbf{validación cruzada out-of-fold} —de modo que cada predicción se obtiene sobre datos no
empleados en el ajuste— y su significancia se contrasta con \textbf{tests de permutación} que
controlan el desbalance de clases. Se prestó especial atención al control de \textit{confounds}
(régimen de la muestra, longitud del informe y artefactos del máximo de atención), validando cada
señal frente a una referencia continua u objetiva para descartar espejismos. La exploración de
señales adicionales se realiza \textbf{sobre los tres conjuntos a la vez}, exigiendo que cualquier
mejora \emph{generalice}; se comprueba que ninguna combinación multi-señal supera de forma
consistente a la probabilidad mínima de contenido por sí sola —las demás señales resultan específicas
de modelo o de dominio—, por lo que la fórmula final privilegia la \textbf{simplicidad robusta} de
una única señal frente a un combinado sobreajustado. Esta se \textbf{calibra} para expresar la
fiabilidad como una probabilidad con su \textbf{intervalo de confianza (\emph{bootstrap})} y bandas
de riesgo. El software empleado es Python (PyTorch, \texttt{transformers}, NumPy) y la cuantización se
realiza con \texttt{bitsandbytes}; el AUC y los intervalos se calculan con estimadores basados en
rangos para garantizar reproducibilidad.
